\section{Block return maps and a common-fibre construction}
\label{sec:block-crt}

The quotient-OR identity for the row-fibred construction must not be
extended to arbitrary cell-dependent blocks. The distinction is already
visible in the order-$12$ counterexample.

\begin{theorem}[Cycle lengths in a Latin block substitution]
\label{thm:block-return}
Let $Q$ be Latin on $A$, and let $F_{ab}$ be Latin on a common set $B$
for every $(a,b)\in A^2$. The operation
\[
 H((a,x),(b,y))=(Q(a,b),F_{ab}(x,y))
\]
is Latin. In each view, same-outer-label line pairs restrict to block
line pairs. For different outer labels, a quotient cycle of length $d$
has a bijective inner return map $\Omega$. Each $e$-cycle of $\Omega$
lifts to one cycle of length $de$. Consequently a view fails exactly
when some block fails in that view, or some odd quotient cycle has an
odd return-map cycle for a lifted line pair.
\end{theorem}

\begin{proof}
In a fixed row, the desired outer symbol determines the outer column;
the block bijection then determines the inner column. The column argument
is identical. For rows $(a,x),(a',x')$, equality of outer symbols sends
$b$ to $b'$ with $Q(a,b)=Q(a',b')$. The inner step is
$F_{a'b'}(x',-)^{-1}F_{ab}(x,-)$. If $a=a'$, then $b=b'$ and this
is exactly a row-pair permutation in one block. Otherwise it projects
to the quotient row-pair map. Exchanging the two inputs proves the
column statement.

For a symbol $(s,z)$, a column $(b,y)$ determines the unique outer row
$a$ with $Q(a,b)=s$; the inner row is obtained by inverting the symbol
line of $F_{ab}$. Equal outer symbols leave $b$ fixed and give a block
symbol-pair permutation. Distinct outer symbols follow the quotient
symbol-pair map. Every block and every pair of its lines occur as these
labels vary, in all three views.

Above a quotient cycle $(b_0,\ldots,b_{d-1})$, write the successive inner
bijections as $h_0,\ldots,h_{d-1}$. A first return to the starting fibre
is $\Omega=h_{d-1}\cdots h_0$. An $e$-cycle of this map closes after
exactly $e$ returns and hence $de$ steps; it cannot close before an
outer return. This accounts for all cycles and proves the parity claim.
\end{proof}

Both frozen order-$12$ examples consist of nine FFF blocks of order $4$
over a TTT quotient of order $3$. Fresh checks verify the return formula
on every physical pair in each view. Thus a failing quotient can be
masked by even inner returns. This neither contradicts the row-fibred
theorem nor the special two-element-kernel equality.

\begin{theorem}[Odd cubefree common-fibre lift]
\label{thm:cubefree-crt}
Let $N>1$ be odd and cubefree, $M=\operatorname{rad}(N)$, and
$e=\operatorname{ord}_M(2)$. If $D>0$ is divisible by $e$ and
$2^D>3(N-1)$, there exists an FFF square of order $N2^D$.
Its two-line cycles have length $2$ or $2h$, where $h$ is the additive
order of a nonzero quotient difference.
\end{theorem}

\begin{proof}
Write $N=\prod_p p^{r_p}$, with $r_p\in\{1,2\}$, and put
$G=\prod_p\mathbb F_p^{r_p}$ and $R=\prod_p\mathbb F_p$.
Choose $B_p(x,y)=xy$ in rank one and
$B_p(x,y)=x_1y_1-\nu_px_2y_2$ in rank two, with $\nu_p$ a nonsquare.
Let $B=(B_p)_p$ and $Q(x)=B(x,x)$. Anisotropy gives
$Q_p(d)\ne0$ exactly when $d_p\ne0$. Hence $Q(d)$ in the additive
group of $R$ and $d$ in $G$ have the same odd order $h(d)$.

The condition on $D$ gives a primitive $M$th root in
$K=\mathbb F_{2^D}$. The additive CRT identification of $R$ with
$\mathbb Z/M\mathbb Z$ gives an injective character $\chi:R\to K^*$.
Define, for a twist array $c:G^2\to K$,
\[
 (x,a)*(y,b)=
 (x+y,\chi(2B(x,y))a+\chi(Q(x))b+c(x,y)).
\]
The two fibre coefficients are nonzero, so every twist gives a Latin
square. The row, column and symbol relative-step formulas are precisely
those in the proof of Theorem~\ref{thm:quadratic-norm-lift}, now evaluated
componentwise in $R$ over a quotient orbit of length $h=h(d)$.

Here is why replacing $p$ by composite $h$ is legitimate. Each slope
exponent is constant or affine in the step index. Its nonzero
coordinates are supported where $d_p\ne0$; each such odd prime divides
$h$, and $\sum_{t=0}^{h-1}t=h(h-1)/2=0$ in that component.
Thus the total slope is one. Composing the later slopes into each of
the two label coefficients gives the same displayed exponents as in
that proof, of the form $A+tQ(d)$ or $A-tQ(d)$. This follows by
expanding the quadratic form in each active component; in an inactive
component $d_p=0$ and the step is constant. The element $\chi(Q(d))$
has exact order $h$, so
\[
 \sum_{t=0}^{h-1}\chi(A\mathbin{\pm}tQ(d))=0.
\]
Both line labels therefore cancel in every view. Every inner return
is $w\mapsto w+\Lambda_C(c)$.

Choose one orientation and starting point per unordered quotient pair
and orbit; conjugate returns have the same nonvanishing condition.
The form $\Lambda_C$ has $2h$ distinct twist cells with nonzero
coefficients. For symbols, the two sets of cells have distinct
quotient symbols; this ensures distinctness just as different quotient
rows or columns do in the other views. A cell lies in one return form
for each other quotient line in each view, a total of $3(N-1)$ forms.
Assign twist cells in a fixed order and enforce a form to be nonzero
when its last cell is assigned. Each closing form forbids one field
value. The field-size inequality leaves an allowed choice at every step.

For different quotient labels the return is now a nonzero translation
of a characteristic-two field, hence has only two-cycles; physical
cycles have length $2h$. For equal quotient labels and distinct fibre
labels $a,b$, the row, column and symbol translations respectively have
nonzero constants
\[
 \chi(2B(x,y)-Q(x))(a+b),\quad
 \chi(Q(x)-2B(x,y))(a+b),\quad
 \chi(-Q(x))(a+b).
\]
They give two-cycles. Every physical pair in all three views is covered.
\end{proof}

For example, $N=11^2\cdot41^2=203401$ has
$\operatorname{ord}_{451}(2)=20$ and $2^{20}>3(N-1)$, giving $h(N)\leq20$.
The product of the two separate prime-square constructions gave the
upper bound $10+20=30$. This improves that particular construction,
not a known minimum; the order-$213281406976$ table is not materialized.
The compact finite records cover $N=9,15,45,75$ and direct return-map
controls. Their checks do not replace the general proof.

\subsection{Limits of the bilinear method}

The finite-field zero-count arguments below are elementary special cases
of the classical Chevalley--Warning method
\cite{Chevalley1935,Warning1935}; their needed details are included.

\begin{proposition}[Scalar quotient boundary]
\label{prop:scalar-boundary}
Let $G$ be a finite odd abelian group and let $\beta:G^2\to K^*$ be
a symmetric bicharacter over a characteristic-two finite field. If
some twist makes
\[
 (x,a)*(y,b)=(x+y,\beta(x,y)^2a+\beta(x,x)b+c(x,y))
\]
row-F, every Sylow subgroup of $G$ is elementary abelian of rank at
most two. Conversely every such group admits an FFF instance over a
sufficiently large field. This is a boundary of this formula only.
\end{proposition}

\begin{proof}
If $d\ne0$ and $\beta(d,d)=1$, composition around a row quotient orbit
of odd length $h$ has slope one. Its two label coefficients are nonzero
scalar multiples of $\sum_{t=0}^{h-1}\beta(d,d)^{\pm t}=1$.
Whatever the twist, one may choose the second row label to cancel the
return translation. This creates an odd physical $h$-cycle. Thus no
nonzero $d$ can have $\beta(d,d)=1$.

An element $u$ of order $p^2$ violates this condition at $d=pu$.
On a rank-three elementary $p$-subgroup, the exponent of
$\beta(d,d)$ is a quadratic form over $\mathbb F_p$ (or is identically
zero). Such a form has a nonzero zero: sum $1-Q(d)^{p-1}$ over
$\mathbb F_p^3$. Each monomial has a variable of exponent below
$p-1$, so its sum is zero modulo $p$. The number of zeros is divisible
by $p$ and includes zero, hence is at least $p$. Necessity follows;
Theorem~\ref{thm:cubefree-crt} proves sufficiency for the stated groups.
\end{proof}

\begin{proposition}[Commuting-matrix fibre bound]
\label{prop:matrix-boundary}
Let $p$ be an odd prime, $r\geq1$ and $D\geq0$.
In the same formula let $G=\mathbb F_p^r$, let the fibre be
$V=\mathbb F_2^D$, and let $\beta$ be a symmetric bicharacter with
commuting image in $\operatorname{GL}(V)$. Row-F requires
$D\geq\operatorname{ord}_p(2)\lceil r/2\rceil$.
If $2^{\operatorname{ord}_p(2)}>3(p^2-1)$, this dimension is attained
by an FFF construction within this ansatz.
\end{proposition}

\begin{proof}
Write $e=\operatorname{ord}_p(2)$. The row-return computation uses only
commutativity: the label coefficients are invertible matrices multiplied
by $S(q)=\sum_{t=0}^{p-1}q^t$ and $S(q^{-1})$, where $q=\beta(d,d)$.
If $q=I$, both sums are $I$, and a choice of the second row label cancels
the return. Thus $\beta(d,d)\ne I$ for every nonzero $d$ is necessary.

The commuting matrices of exponent $p$ diagonalize simultaneously over
$\mathbb F_{2^e}$. Each nontrivial eigencharacter has a Frobenius orbit
of length $e$, so there are at most $t=\lfloor D/e\rfloor$ such orbits.
One representative from each yields a quadratic exponent form $Q_i$ on
$\mathbb F_p^r$. If $r>2t$, summing
$\prod_i(1-Q_i(d)^{p-1})$ shows, by the same degree argument as above,
that their common zero count is divisible by $p$. A nonzero common
zero would give $\beta(d,d)=I$, a contradiction. Hence $r\leq2t$.
For attainment, split $G$ into rank-two summands and at most one rank-one
summand. Apply Theorem~\ref{thm:cubefree-crt} to each over
$\mathbb F_{2^e}$ and take their direct product. The resulting block
diagonal bicharacter has exactly the required fibre dimension.
\end{proof}

These propositions concern the exact scalar or commuting-matrix
formula, not nonlinear blocks, arbitrary quotient operations, or the
unrestricted minimum $h(m)$. They provide no new exclusion of order $18$.
In particular, having a fixed vector is not the same as being the identity
matrix and cannot replace the condition used in the last proof.
